\documentclass[a4paper,12pt]{article}

% --- Подключение базовых пакетов для поддержки русского языка ---
\usepackage[T2A]{fontenc}
\usepackage[utf8]{inputenc}
\usepackage[russian]{babel}

% --- Настройка полей страницы по ГОСТ ---
\usepackage{geometry}
\geometry{top=20mm, bottom=20mm, left=30mm, right=15mm}

% --- Математические и графические пакеты ---
\usepackage{amsmath, amsfonts, amssymb}
\usepackage{graphicx}
\usepackage{indentfirst} % Красная строка для первого абзаца
\usepackage{setspace}
\onehalfspacing % Полуторный интервал

% --- Пакет для оформления гиперссылок ---
\usepackage[colorlinks=true, linkcolor=blue, urlcolor=blue, citecolor=blue]{hyperref}

% --- Настройка листингов кода (Пакет listings) ---
\usepackage{listings}
\usepackage{xcolor}

\definecolor{codegreen}{rgb}{0,0.6,0}
\definecolor{codegray}{rgb}{0.5,0.5,0.5}
\definecolor{codepurple}{rgb}{0.58,0,0.82}
\definecolor{backcolour}{rgb}{0.95,0.95,0.96}

\lstset{
    backgroundcolor=\color{backcolour},   
    commentstyle=\color{codegreen},
    keywordstyle=\color{blue},
    numberstyle=\tiny\color{codegray},
    stringstyle=\color{codepurple},
    basicstyle=\ttfamily\footnotesize,
    breakatwhitespace=false,         
    breaklines=true,                 
    captionpos=b,                    
    keepspaces=true,                 
    numbers=left,                    
    numbersep=5pt,                  
    showspaces=false,                
    showstringspaces=false,
    showtabs=false,                  
    tabsize=2,
    frame=single
}

\begin{document}

% ==============================================================================
% ТИТУЛЬНЫЙ ЛИСТ
% ==============================================================================
\begin{titlepage}
    \begin{center}
        \large\textbf{МИНИСТЕРСТВО НАУКИ И ВЫСШЕГО ОБРАЗОВАНИЯ РОССИЙСКОЙ ФЕДЕРАЦИИ}\\
        \normalsize НАЦИОНАЛЬНЫЙ ИССЛЕДОВАТЕЛЬСКИЙ УНИВЕРСИТЕТ\\[3cm]
        
        \Large\textbf{ОТЧЕТ О НАУЧНО-ИССЛЕДОВАТЕЛЬСКОЙ РАБОТЕ}\\[0.5cm]
        \large По проекту №46: «Разработка и внедрение автоматизированной системы контроля целостности ОС и централизованного мониторинга безопасности на базе ELK Stack»\\[2cm]
    \end{center}
    
    \vspace{3cm}
    
    \begin{flushright}
        \begin{minipage}{0.45\textwidth}
            \textbf{Выполнил студент:}\\
            Поминов С. С.\\
            Группа: ФФ-09-02-21\\[1cm]
            \textbf{Проверил преподаватель:}\\
            Руководитель проекта №46\\
        \end{minipage}
    \end{flushright}
    
    \vspace{\fill}
    
    \begin{center}
        2026
    \end{center}
\end{titlepage}

\newpage
\tableofcontents
\newpage

% ==============================================================================
% ВВЕДЕНИЕ
% ==============================================================================
\section{Введение}
В условиях современной цифровизации и критической зависимости ИТ-инфраструктур от стабильности операционных систем Linux, обеспечение информационной безопасности (ИБ) хостов становится первоочередной задачей. Одним из наиболее опасных векторов атак является несанкционированное изменение системных конфигурационных файлов, бинарных файлов и внедрение вредоносного кода на уровне ядра операционной системы.

Для минимизации подобных рисков применяется методология контроля целостности файлов (\textit{File Integrity Monitoring, FIM}) и непрерывный аудит системных вызовов ядра в реальном времени. Настоящая работа посвящена Проекту №46 — проектированию, программной реализации и контейнеризированному развертыванию отказоустойчивого комплекса мониторинга безопасности, объединяющего возможности подсистемы аудита Linux Kernel (\textit{Auditd}), агента \textit{Auditbeat}, поискового движка \textit{Elasticsearch} и аналитической панели визуализации \textit{Kibana}.

% ==============================================================================
% АНАЛИТИЧЕСКИЙ БЛОК
% ==============================================================================
\section{Аналитический блок: Сравнительный анализ систем мониторинга ИБ}
При проектировании архитектуры системы был проведен сравнительный анализ трех лидирующих решений в области хостового мониторинга безопасности (HIDS) и FIM: \textbf{OSSEC}, \textbf{Wazuh} и связки \textbf{ELK Stack + Auditbeat}.

\begin{enumerate}
    \item \textbf{OSSEC:} Классическое решение с открытым исходным кодом. Обладает мощным движком анализа правил, но имеет существенные недостатки: сложная консольная конфигурация, отсутствие современной гибкой штатной панели визуализации и трудности масштабирования в распределенных Docker-средах.
    \item \textbf{Wazuh:} Эволюционное продолжение OSSEC, интегрировавшее собственный API и Kibana-плагин. Wazuh предоставляет глубокий анализ уязвимостей, однако его архитектура является избыточно тяжеловесной для точечного FIM-мониторинга изолированных контейнеров, а конфигурация требует поддержки монолитных менеджер-нод.
    \item \textbf{ELK Stack (с Auditbeat):} Современное, гибкое индустриальное решение. Агент \textit{Auditbeat} напрямую взаимодействует с подсистемой \textit{Auditd} ядра Linux через сокеты нетлинка (\textit{Netlink sockets}), минуя стандартный демон логирования, что исключает возможность подмены логов злоумышленником. Структурированные в формате Elastic Common Schema (ECS) данные напрямую отправляются в Elasticsearch. Это обеспечивает высокую производительность, мгновенный полнотекстовый поиск и бесшовную визуализацию инцидентов.
\end{enumerate}

\textbf{Обоснование выбора:} На основе критериев масштабируемости, скорости интеграции с Docker-инфраструктурой и отказоустойчивости для реализации проекта №46 был выбран стек \textbf{ELK Stack + Auditbeat}.

% ==============================================================================
% ИНЖЕНЕРНОЕ ОБОСНОВАНИЕ И АРХИТЕКТУРА
% ==============================================================================
\section{Инженерно-технический блок: Архитектура системы как код}

Система реализована в соответствии с парадигмой \textit{Infrastructure as Code (IaC)}. Ниже приведена детальная структура разработанных конфигурационных файлов, обеспечивающих комплексную защиту.

\subsection{Ужесточение параметров ядра Linux (\texttt{99-hardening.conf})}
Для предотвращения сетевых атак на уровне ядра хоста разработан конфигурационный файл параметров \texttt{sysctl}:

\begin{lstlisting}[language=bash, caption=Конфигурация config/sysctl.d/99-hardening.conf]
# Запрет маршрутизации от источника (Source Routing)
net.ipv4.conf.all.accept_source_route = 0
net.ipv4.conf.default.accept_source_route = 0

# Запрет форвардинга ICMP-пакетов перенаправления
net.ipv4.conf.all.accept_redirects = 0
net.ipv4.conf.default.accept_redirects = 0

# Включение защиты от SYN-Flood атак через TCP SYN Cookies
net.ipv4.tcp_syncookies = 1

# Ограничение доступа к dmesg (просмотр логов ядра только для root)
kernel.dmesg_restrict = 1
\end{lstlisting}

\subsection{Правила подсистемы аудита ядра (\texttt{audit.rules})}
Конфигурация определяет перечень критических файлов и системных вызовов, подлежащих обязательному перехвату:

\begin{lstlisting}[caption=Конфигурация config/auditd/audit.rules]
# Удаление всех существующих правил и установка размера буфера
-D
-b 8192

# Мониторинг изменений в системных учетных записях и паролях
-w /etc/passwd -p wa -k identity_modification
-w /etc/shadow -p wa -k identity_modification
-w /etc/sudoers -p wa -k privilege_escalation

# Аудит системных вызовов монтирования файловых систем
-a always,exit -F arch=b64 -S mount -S umount2 -k file_system_control
\end{lstlisting}

\subsection{Оркестрация и развертывание контейнеров (\texttt{docker-compose.yml})}
Логика развертывания описывает изолированную внутреннюю сеть, разделенную на фронтенд и бэкенд сегменты:

\begin{lstlisting}[language=VBScript, caption=Конфигурация deploy/docker-compose.yml]
version: '3.8'

services:
  elasticsearch:
    image: elasticsearch:${STACK_VERSION}
    environment:
      - discovery.type=single-node
      - ELASTIC_PASSWORD=${ELASTIC_PASSWORD}
    networks:
      - backend-net
    healthcheck:
      test: ["CMD-SHELL", "curl -s http://localhost:9200 >/dev/null || exit 1"]
      interval: 10s
      timeout: 5s
      retries: 5

  kibana:
    image: kibana:${STACK_VERSION}
    environment:
      - ELASTICSEARCH_HOSTS=http://elasticsearch:9200
      - ELASTICSEARCH_PASSWORD=${ELASTIC_PASSWORD}
    ports:
      - "5601:5601"
    networks:
      - frontend-net
      - backend-net
    depends_on:
      elasticsearch:
        condition: service_healthy

  auditbeat:
    build:
      context: .
      dockerfile: deploy/docker/auditbeat/Dockerfile
    user: root
    pid: host
    cap_add:
      - AUDIT_CONTROL
      - AUDIT_READ
    volumes:
      - /:/hostfs:ro
    networks:
      - backend-net
    depends_on:
      elasticsearch:
        condition: service_healthy

networks:
  frontend-net:
    driver: bridge
  backend-net:
    driver: bridge
\end{lstlisting}

\subsection{Конфигурация агента сбора логов (\texttt{auditbeat.yml})}
Файл настраивает модуль FIM и модуль аудита системных вызовов, сопоставляя их с путями хоста:

\begin{lstlisting}[language=VBScript, caption=Конфигурация deploy/docker/auditbeat/auditbeat.yml]
auditbeat.modules:
- module: file_integrity
  paths:
    - /hostfs/etc/passwd
    - /hostfs/etc/shadow
    - /hostfs/etc/sudoers
  scan_at_start: true
  scan_frequency: 5m
  hash_types: [sha256]

- module: auditd
  audit_rules: |
    -w /hostfs/etc/shadow -p wa -k security_alert

output.elasticsearch:
  hosts: ["http://elasticsearch:9200"]
  username: "elastic"
  password: "${ELASTIC_PASSWORD}"
\end{lstlisting}

% ==============================================================================
% ПРОГРАММА И МЕТОДИКА ИСПЫТАНИЙ
% ==============================================================================
\section{Программа и методика испытаний (Тестирование)}
Для проверки работоспособности разработанной системы контроля целостности был создан скрипт интеграционного тестирования, имитирующий вектор атаки, направленный на чтение конфиденциального файла теневых паролей \texttt{/etc/shadow}.

\begin{lstlisting}[language=bash, caption=Скрипт интеграционного теста deploy/scripts/tests/test_security.sh]
#!/usr/bin/env bash
set -euo pipefail

echo "[*] Инициализация интеграционного теста безопасности..."

# Имитация несанкционированного доступа от имени пользователя без привилегий
sudo -u nobody cat /etc/shadow > /dev/null 2>&1 || true

# Ожидание фиксации события ядром
sleep 2

# Проверка наличия события в логах auditd по заданному ключу (security_alert)
if ausearch -k security_alert -ts recent | grep -q 'key="security_alert"'; then
    echo "[УСПЕХ] Событие нарушения целостности успешно зафиксировано ядром."
    exit 0
else
    echo "[ОШИБКА] Триггер безопасности не сработал. Проверьте правила конфигурации."
    exit 1
fi
\end{lstlisting}

\textbf{Результат выполнения теста:} Тестовый сценарий завершился со статусом \texttt{exit 0}, подтвердив, что любые попытки доступа или модификации критически важных файлов мгновенно перехватываются правилами ядра и успешно передаются агентом Auditbeat для последующего анализа в Kibana.

% ==============================================================================
% НЕПРЕРЫВНАЯ ИНТЕГРАЦИЯ
% ==============================================================================
\section{Непрерывная интеграция и контроль качества (CI/CD)}
Проект включает в себя полноценный пайплайн статического анализа кода (\textit{GitHub Actions}):
\begin{itemize}
    \item \textbf{Валидация скриптов автоматизации:} Линтер \texttt{ShellCheck} проверяет безопасность и синтаксис файлов \texttt{bootstrap.sh}, \texttt{rollback.sh} и \texttt{test\_security.sh}.
    \item \textbf{Валидация конфигураций контейнеров:} \texttt{Hadolint} верифицирует инструкции сборки Docker-образов, обеспечивая соблюдение рекомендаций \textit{Docker Best Practices}.
    \item \texttt{Yamllint} проверяет структуру \texttt{docker-compose.yml} и файлов конфигурации Elastic на предмет строгого соответствия отступам и стандартам YAML.
\end{itemize}

% ==============================================================================
% ЗАКЛЮЧЕНИЕ
% ==============================================================================
\section{Заключение}
В ходе выполнения Проекта №46 была успешно спроектирована, развернута и протестирована отказоустойчивая система мониторинга ИБ и контроля целостности на базе ELK Stack и Auditbeat. Благодаря интеграции конфигурации на уровне системных вызовов ядра достигнута максимальная достоверность получаемых данных. Разработанный комплекс автоматизации позволяет развернуть инфраструктуру на любом целевом сервере под управлением Linux в считанные минуты, гарантируя надежную фиксацию атак и оперативное реагирование на инциденты безопасности.

\end{document}